\subsection{Phylogeny mostly gets it right: where IBA fails and why}\label{sec:dd-iba}
\textbf{Claim.} IBA is the best-performing inferred evidence (\accIBA\% accepted, \rejIBA\% rejected,
vs ISO \rejISO\%), because PAINT curators place functions at tree nodes; its failures are
structured and fixable at the node.
\textbf{Data.} \proj{IBA\_REVIEW.md}: 16 recurring failure patterns, 79 worked cases;
\texttt{propagation\_review} blocks on 691 reviews. \proj{PANTHER\_IBA\_REVIEW.md}: all 160 IBAs on
41 \textit{S.\ pombe} genes rebuilt from source node; 148 kept. Both directions: IBA too
conservative for core MF in human (\texttt{reports/iba-too-conservative-core-mf.md}).
\textbf{Worked examples.} \todo{one pseudoenzyme inheriting catalysis; one neo-functionalised
subfamily; one cross-kingdom extension.} Conserved-function recovery: flagellar/T3SS export
ATPases (FliI/SctN) annotated as F-type ATP synthase through PAINT, TreeGrafter, InterPro2GO and
GO logical inference (\proj{CONSERVED\_FUNCTION\_RECOVERY.md}).
\textbf{Message.} Do not read donor count as evidential weight; argue with node placement.
